% SOURCE_AUTHORITY: sga5_fr_workpass.tex, lines 13660--13693 (LF-normalized unit).
% SOURCE_SHA256: 791F4EFFC5E02832D5D77ED03518C8156D6F07E4C8238B03545DB93D883FBB28
\medskip

\begin{prooftext}[Demonstração da proposição 4.6]
Já se observou que podemos reduzir ao caso $g = e$, $\varphi = \id_G$. Se denotarmos por $G_1$ o núcleo do homomorfismo
\[
G \longrightarrow \Aut_k(\fm/\fm^2) \quad ,
\]
então, graças à teoria da ramificação, tem-se:
\[
\card(G_1) = p^t \quad .
\]

Por conseguinte, com as notações do lema 4.7, tem-se:
\[
\card(G) = p^t c \quad .
\]

Como, por outro lado, $c^{\id_G}(e) = \card(G)$, o enunciado da proposição 4.6 no caso $g = e$, $\varphi = \id_G$ e, por conseguinte, em geral, resulta do lema 4.7.
\end{prooftext}

\medskip
\noindent
\textbf{Corolário 4.8.} \emph{Para cada $\ell \neq \car(k)$, a imagem do número racional
\[
\frac{(\nu(g_R^{-1}f'_R) - 1)}{c^\varphi(g)}
\]
pela inclusão canônica
\[
\bQ \hookrightarrow \bQ_\ell
\]
é um inteiro $\ell$-ádico.}

\medskip
\noindent
\textbf{Proposição 4.9.} \emph{A conjectura 4.5 é verdadeira quando a dimensão do esquema $X$ é igual a 1.}
